\section*{Chapter \ref{ch:hf:lead-lag-and-cross-venue-trading} --- Lead--Lag and Cross-Venue Trading}

\begin{solution}{exo:hf:lead-lag-and-cross-venue-trading:1}
Market 2 corrects (large $\alpha_2$), so market 1 leads; $\gamma_1=0.18/(0.18+0.02)=0.9$.
\end{solution}

\begin{solution}{exo:hf:lead-lag-and-cross-venue-trading:2}
$67.5$ trades a session $\times0.363\times100\times\$0.01=\$24.50$.
\end{solution}

\begin{solution}{exo:hf:lead-lag-and-cross-venue-trading:3}
$97/7\approx14$.
\end{solution}

\begin{solution}{exo:hf:lead-lag-and-cross-venue-trading:4}
When the lead is shorter than the grid step, both markets' reactions to a piece of news fall in the same interval: their innovations are
correlated and the Cholesky ordering decides who gets the shared part. When the lead spans several steps, the leader's innovation comes first and
the follower's is mostly the error correction, so the innovations are nearly uncorrelated and the bounds meet.
\end{solution}

\begin{solution}{exo:hf:lead-lag-and-cross-venue-trading:5}
Once the follower has moved, the trade is taken at the new price with no gap left to close: the expected move is zero and the taker still pays half
the spread and the fee, $-0.6$ ticks, close to the $-0.54$ measured at one second.
\end{solution}

\begin{solution}{exo:hf:lead-lag-and-cross-venue-trading:6}
Between 300 milliseconds ($+0.24$) and 500 ($-0.22$).
\end{solution}

\begin{solution}{exo:hf:lead-lag-and-cross-venue-trading:7}
0.444 ticks a share at one millisecond and 0.142 at 100: with a lead of 0.2 seconds, 100 milliseconds already costs two thirds of the edge.
\end{solution}

\begin{solution}{exo:hf:lead-lag-and-cross-venue-trading:8}
The share depends on the sampling step: on one-second data a lead of milliseconds is invisible and the shares mostly reflect which market's noise is
smaller. The fund may matter at the frequencies where the desk trades; measure on its own ticks and at its own horizon.
\end{solution}

\begin{solution}{pb:hf:lead-lag-and-cross-venue-trading:1}
\begin{enumerate}
\item Most price discovery in the E-mini for the S\&P 500 and Nasdaq-100; shared between the regular future and the fund for the S\&P 400; the S\&P 500
fund leads sector funds far more than the reverse.
\item Futures carried most price discovery for five- and ten-year notes from 1999 on.
\item Median correlation 0.10 at 10 ms and 0.008 at 1 ms in 2011.
\item See \cref{def:hf:lead-lag-and-cross-venue-trading:lead}.
\item $\Delta p_t=\alpha(p^1_{t-1}-p^2_{t-1})+\sum\Gamma_i\Delta p_{t-i}+e_t$; $\gamma_1=\alpha_2/(\alpha_2-\alpha_1)$; the information share from the Cholesky
factor, in both orderings.
\item Information share 0.52--0.73, 0.81--0.90, 0.97--0.98 and 0.998 for leads of 0.05, 0.2, 0.5 and 1 second; component share 0.57, 0.73, 0.87,
0.98; the bounds narrow as the grid resolves the lead.
\item Sampling on a grid adds the Epps bias and blurs leads shorter than the step.
\item 0.05 and 0.25; 0.225 and 0.25; 0.475 and 0.575; 1.075 and 1.1 seconds.
\item The default books followed their efficient price over seconds, so the cross-correlation peak was flat; the estimates were $-1.65$ and $+1.0$
seconds.
\item Cross the follower's spread for one lot when the leader has moved at least a tick more over the last second; marked at the follower's mid five
seconds later, after half the spread and a 0.1-tick fee.
\item 0.485, 0.43, 0.24, $-0.22$ and $-0.54$ ticks a share.
\item The planted lead is half a second; a real future--fund lead was a few milliseconds in 2011, so the curve would be compressed into them.
\item It withdraws or skews its stale side when the leader moves (chapter 7).
\item Planted 0.5 s, estimated 0.475 and 0.575 s; the leader's information share 0.97--0.98; the edge disappears between 300 and 500 ms.
\item Liquidity, hours, fees, outages and competition move price discovery; trading the lead itself shortens it.
\item Leading crypto venue to lagging venue: Makarov and Schoar found gaps much larger across countries, where capital controls slow arbitrage.
\item Rolling estimates of the lead and the information shares by time of day, with alerts when they move beyond their usual range.
\item What the taker earns is the move it trades ahead of; once faster than the lead, being faster still adds little in this market (0.485 at one
millisecond against 0.463 at fifty).
\item Nothing is earned by the unused speed while its fixed cost is paid (chapter 1's operating leverage).
\item Speed relative to the others trading the same lead.
\end{enumerate}
\end{solution}

\begin{solution}{iq:hf:lead-lag-and-cross-venue-trading:1}
The move may already be priced in the fund's other venues; the fund's market makers may have withdrawn and the spread widened; the future's move
may be a single print that reverts; the gap may not cover the spread and fees at the time the order arrives.
\iqlookfor{costs, stale data, and the chance of being late.}
\end{solution}

\begin{solution}{iq:hf:lead-lag-and-cross-venue-trading:2}
Grid sampling pairs increments over intervals in which one series may not have moved yet, so part of the common move is missed (the Epps effect).
Hayashi--Yoshida sums products of every pair of overlapping increments on the series' own ticks, so no move is dropped.
\iqlookfor{the source of the bias and the overlap idea.}
\end{solution}

\begin{solution}{iq:hf:lead-lag-and-cross-venue-trading:3}
The innovations of the two markets are correlated, and the Cholesky decomposition assigns the correlated part to whichever market is ordered first;
the two orderings bound the share. They are close when the innovations are nearly uncorrelated: a fine grid, or a lead longer than the step.
\iqlookfor{the ordering argument.}
\end{solution}

\begin{solution}{iq:hf:lead-lag-and-cross-venue-trading:4}
Timestamp at the network card with hardware stamps disciplined by a common time source (precision time protocol to a satellite reference), record the
exchanges' own timestamps too, and measure the path delay between the sites (One Quant Book 14, chapter 4).
\iqlookfor{hardware timestamps and a shared clock.}
\end{solution}

\begin{solution}{iq:hf:lead-lag-and-cross-venue-trading:5}
Someone faster trades the same gaps: your fills come later in each opportunity. Check your latency rank (fill times against the leader's move), the
follower's quote-withdrawal speed and the lead's own length.
\iqlookfor{competition and how to measure it.}
\end{solution}

\begin{solution}{iq:hf:lead-lag-and-cross-venue-trading:6}
$\alpha_\perp=(\alpha_2,0)$, so $\gamma=(1,0)$. With market 1 ordered first its information share is one; ordered second it is $1-\rho^2$, with $\rho$
the correlation of the two innovations: the bounds are $[1-\rho^2,1]$.
\iqlookfor{the algebra and the role of innovation correlation.}
\end{solution}
